\section{Local spectral theory and global-to-local convergence}
\label{sec:local}

Global energy decay eventually places the iterates near the optimizing
Gaussian, but it does not determine the local rate.  After optimizer
whitening, the derivative of the Fisher--Rao field contains first- and
second-Hermite correlations of the target Hessian.  This section identifies
the corresponding self-adjoint generator, bounds its spectrum, and then
controls the nonlinear dynamics by retaining the Gaussian part of the field
exactly.

\subsection{Optimizer whitening and the score generator}
\label{sec:local-linearization}

Work in the optimizer-whitened coordinates introduced in
\Cref{sec:introduction-bottlenecks}.  Thus \(a_\star=(0,\Id)\),
\[
  \E_{\N(0,\Id)}\nabla\widehat V(Z)=0,\qquad
  \E_{\N(0,\Id)}\nabla^2\widehat V(Z)=\Id,
\]
and
\[
  \alphastar\Id\preceq\nabla^2\widehat V(z)
  \preceq\betastar\Id.
\]
For \(\xi=(u,X)\in\R^d\times\Sym(d)\), use the equilibrium Fisher--Rao norm
\begin{equation}
  \norm{\xi}_\star^2
  :=\norm u^2+\frac12\norm X_{\Frob}^2.
  \label{eq:local-equilibrium-norm}
\end{equation}
Let \(B(Z):=\nabla^2\widehat V(Z)\).  The first- and second-Hermite score
operators are
\begin{align}
  \mathcal T[X]_k
  &:=\E\!\left[\Tr(XB(Z))Z_k\right],
  &
  \mathcal T^\ast[u]_{ij}
  &:=\E\!\left[(u^\top Z)B_{ij}(Z)\right],
  \label{eq:local-first-hermite}\\
  \mathcal H[X]_{ij}
  &:=\E\!\left[
      B_{ij}(Z)\bigl(Z^\top XZ-\Tr X\bigr)
    \right].
  \label{eq:local-second-hermite}
\end{align}
We use the operator norm
\begin{equation}
  \norm{\mathcal T}
  :=
  \sup_{X\neq0}
  \frac{\norm{\mathcal T[X]}}{\norm X_{\Frob}}.
  \label{eq:local-first-hermite-operator-norm}
\end{equation}
Gaussian integration by parts gives
\[
  u^\top\mathcal T[X]=\Tr(\mathcal T^\ast[u]X),
  \qquad
  \Tr(\mathcal H[X]Y)=\Tr(X\mathcal H[Y]).
\]
Consequently, the negative Jacobian of the Fisher--Rao field at equilibrium
is the self-adjoint operator
\begin{equation}
  \mathscr L_\star(u,X)
  =
  \left(
    u+\frac12\mathcal T[X],
    X+\mathcal T^\ast[u]+\frac12\mathcal H[X]
  \right).
  \label{eq:local-generator}
\end{equation}
Write
\[
  \gstar:=\lambda_{\min}(\mathscr L_\star),
  \qquad
  \Lstar:=\lambda_{\max}(\mathscr L_\star).
\]

The dimension dependence enters through a centered diagonal block.  Set
\begin{equation}
  t_0:=\max\{1,\log(4\betastar)\},
  \qquad
  \Gamma:=
  \min\!\left\{
    \betastar-1,\,
    \sqrt d\bigl(4\sqrt{2t_0}+4t_0+4\bigr)
  \right\}.
  \label{eq:local-Gamma}
\end{equation}
In particular,
\(\Gamma=O(1+\sqrt d\log\kappastar)\).

\begin{theorem}[Spectral sandwich for the local generator]
\label{thm:local-spectral-sandwich}
Under the whitened curvature bounds,
\begin{equation}
  \boxed{
  \max\!\left\{\alphastar,\frac1{4+\Gamma}\right\}
  \leq\gstar\leq\Lstar
  \leq
  \min\!\left\{\betastar,\frac{4+\Gamma}{2}\right\}.}
  \label{eq:local-spectral-sandwich}
\end{equation}
If the target is Gaussian, then
\(\mathcal T=\mathcal H=0\) and
\(\gstar=\Lstar=1\).
\end{theorem}

\begin{proof}
Diagonalize \(X=\diag(\lambda)\) and define
\[
  \widetilde T_{ki}:=\E[B_{ii}(Z)Z_k],
  \qquad
  D_{ij}:=\E[B_{ii}(Z)(Z_j^2-1)].
\]
The restriction of \(\mathcal H\) to diagonal matrices is self-adjoint, so
\(D\) is symmetric.  For
\(\eta(Z)=u+\diag(\lambda)Z\), two Gaussian integrations by parts give
\begin{equation}
  \E[\eta^\top B\eta]
  =
  \norm u^2+2u^\top\widetilde T\lambda
  +\norm\lambda^2+\lambda^\top D\lambda.
  \label{eq:local-schur-identity}
\end{equation}
Positivity of \(B\) implies
\(\widetilde T^\top\widetilde T\preceq\Id+D\).
The pointwise upper curvature bound gives
\(\lambda_{\max}(D)\leq\betastar-1\).

For the second branch of \(\Gamma\), normalize \(\norm\lambda=1\) and put
\(Q_X(Z)=Z^\top XZ-\Tr X\).  The exact representation
\[
  \lambda^\top D\lambda
  =\E\!\left[\Tr(XB(Z))Q_X(Z)\right]
\]
is combined with
\(\E\Tr(|X|B)=\Tr|X|\leq\sqrt d\) and a two-sided Gaussian
quadratic-form tail bound.  Splitting at \(t_0\) yields
\[
  \abs{\lambda^\top D\lambda}
  \leq
  \sqrt d\bigl(4\sqrt{2t_0}+4t_0+4\bigr).
\]
The quadratic form of \(\mathscr L_\star\) is
\[
  Q(u,\diag\lambda)
  =
  \norm u^2+u^\top\widetilde T\lambda
  +\frac12\norm\lambda^2+\frac14\lambda^\top D\lambda.
\]
Completing the square in \(u\), using the Schur bound, and taking
\(\nu=(4+\Gamma)^{-1}\) gives
\(Q\geq\nu\norm{(u,X)}_\star^2\).
The same estimates give
\(Q\leq(4+\Gamma)\norm{(u,X)}_\star^2/2\).
Finally, applying \eqref{eq:local-schur-identity} with \(u\) replaced by
\(2u\) gives the exact representation
\[
  Q(u,X)
  =\frac14\E\!\left[(2u+XZ)^\top B(Z)(2u+XZ)\right]
   +\frac14\norm X_{\Frob}^2.
\]
Since \(\E\norm{2u+XZ}^2=4\norm u^2+\norm X_{\Frob}^2\) and
\(\alphastar\leq1\leq\betastar\), the pointwise curvature sandwich gives
the direct bounds
\(\alphastar\norm{\xi}_\star^2\leq Q(\xi)
\leq\betastar\norm{\xi}_\star^2\).
Taking the stronger lower and upper estimates proves
\eqref{eq:local-spectral-sandwich}.
\end{proof}

\begin{corollary}[Linearized deterministic maps]
\label{cor:local-linear-maps}
Let \(T_{\Riem,h}\) and \(T_{\KL,h}\) denote the Riemannian and KL maps of
\Cref{sec:geometry-fr-algorithms}, expressed in whitened coordinates.  If
\(0<h\leq2/(4+\Gamma)\), then
\begin{align}
  \norm{DT_{\Riem,h}(a_\star)\xi}_\star
  &\leq
  \left(1-\frac{h}{4+\Gamma}\right)\norm{\xi}_\star,
  \label{eq:local-linear-riem}\\
  \norm{DT_{\KL,h}(a_\star)\xi}_{\star,h}
  &\leq
  \left(1-\frac{h}{4+2h+\Gamma}\right)
  \norm{\xi}_{\star,h},
  \label{eq:local-linear-kl}
\end{align}
where
\(\norm{(u,X)}_{\star,h}^2
=\norm u^2+(1+h)\norm X_{\Frob}^2/2\).
\end{corollary}

\subsection{The Gaussian core and a nonlinear local region}
\label{sec:local-gaussian-core}

Write the whitened Fisher--Rao field as
\[
  F=\FGauss+H,
  \qquad
  \FGauss(m,C)=(-Cm,C-C^2).
\]
The field \(\FGauss\) is the exact natural gradient for the standard Gaussian
target.  It is not treated as a quadratic Taylor remainder.  With
\(\xi=(m,C-\Id)\),
\begin{equation}
  \inner{\xi}{\FGauss(m,C)}_\star
  =
  -m^\top Cm-\frac12\Tr\!\left(C(C-\Id)^2\right).
  \label{eq:local-gaussian-dissipation}
\end{equation}
Thus its dissipation remains useful well beyond the infinitesimal regime.

To turn an energy level into a covariance band, define
\[
  \phi_{\alphastar}(s)
  :=s-1-\log s+\frac{\alphastar}{2}(s-1)^2.
\]
If \(s_1,\ldots,s_d\) are the eigenvalues of \(C^{1/2}\), strong convexity
and the equilibrium first-order conditions imply
\begin{equation}
  \Delta(m,C)
  \geq
  \frac{\alphastar}{2}\norm m^2+
  \sum_{i=1}^d\phi_{\alphastar}(s_i).
  \label{eq:local-entropy-coercivity}
\end{equation}
Let \(s_-(\delta)<1<s_+(\delta)\) be the two roots of
\(\phi_{\alphastar}(s)=\delta\), and set
\[
  \ell_\delta=s_-(\delta)^2,\qquad
  U_\delta=s_+(\delta)^2.
\]
Then the sublevel
\(\mathcal U_\delta:=\{a:\Delta(a)\leq\delta\}\) satisfies
\(\ell_\delta\Id\preceq C\preceq U_\delta\Id\).
Define the explicit norm-conversion factors
\begin{align}
  \vartheta(\delta)
  &:={}
  \sup_{\substack{s\in[s_-(\delta),s_+(\delta)]\\s\neq1}}
  \frac{(s^2-1)^2}{2\phi_{\alphastar}(s)}
  \leq
  \frac{(1+s_+(\delta))^2}
       {\alphastar+s_+(\delta)^{-2}},
  \label{eq:local-vartheta}\\
  \mathfrak b(\delta)^2
  &:={}
  \max\left\{\frac2{\alphastar},\vartheta(\delta)\right\},
  &
  \mathfrak b_h(\delta)^2
  &:={}
  \max\left\{\frac2{\alphastar},(1+h)\vartheta(\delta)\right\}.
  \label{eq:local-b-factors}
\end{align}
Then
\begin{equation}
  a\in\mathcal U_\delta
  \quad\Longrightarrow\quad
  \begin{aligned}
    \norm{a-a_\star}_\star
    &\leq\mathfrak b(\delta)\sqrt\delta,\\
    \norm{a-a_\star}_{\star,h}
    &\leq\mathfrak b_h(\delta)\sqrt\delta.
  \end{aligned}
  \label{eq:local-norm-conversion}
\end{equation}

Let \(\mathcal H_\delta\) be the star-shaped hull of
\(\mathcal U_\delta\) about \(a_\star\).  Every point in this compact hull
obeys the same covariance band and the same bound
\eqref{eq:local-norm-conversion}: its covariance is a convex combination of
\(\Id\) and a covariance in \(\mathcal U_\delta\), while its displacement
from \(a_\star\) is shortened by the same interpolation factor.  Define the
non-Gaussian second-differential modulus
\begin{equation}
  \Kng(\delta)
  :=
  \sup_{\substack{a\in\mathcal H_\delta\\\zeta\neq0}}
  \frac{\norm{D^2H(a)[\zeta,\zeta]}_\star}
       {\norm{\zeta}_\star^2}.
  \label{eq:local-nongaussian-modulus}
\end{equation}
To make this modulus quantitative, put
\begin{align}
  \mathcal B_\delta
  &:={}
  \sup_{a\in\mathcal H_\delta}
  \norm{A(a)-\Id}_{\op},
  \label{eq:local-Bdelta}\\
  \Sigma_\delta
  &:={}
  \sup_{a\in\mathcal H_\delta}
  \left(
    \E_{\rho_a}
    \norm{\nabla^2\widehat V(X)-A(a)}_{\Frob}^2
  \right)^{1/2},
  &
  \overline\Sigma_\delta
  &:={}
  \min\{\Sigma_\delta,\betastar-\alphastar\}.
  \label{eq:local-Sigmadelta}
\end{align}
The Gaussian score calculation in \Cref{app:local-moduli} gives the fully
explicit estimate
\begin{align}
  \Kng(\delta)
  &\leq
  \left\{
    K_{\mathrm{ng},m}(\delta)^2
    +\frac12K_{\mathrm{ng},C}(\delta)^2
  \right\}^{1/2},
  \label{eq:local-Kng-explicit}\\
  K_{\mathrm{ng},m}(\delta)
  &:={}
  \frac{U_\delta^{3/2}}{\ell_\delta^2}
  \left(
    \sqrt3\,\overline\Sigma_\delta+\Sigma_\delta
    +\sqrt{\Sigma_\delta^2+2\mathcal B_\delta^2}
  \right),
  \label{eq:local-Kng-mean}\\
  K_{\mathrm{ng},C}(\delta)
  &:={}
  \frac{U_\delta^2}{\ell_\delta^2}
  \left((4\sqrt2+\sqrt6)\Sigma_\delta+4\mathcal B_\delta\right).
  \label{eq:local-Kng-covariance}
\end{align}
Thus \(\Kng(\delta)\) measures only departure from the Gaussian core: all
three non-Gaussian moduli vanish identically for a Gaussian target.  Bounded
Hessian and Gaussian score differentiation suffice for finiteness; classical
third and fourth derivatives of \(V\) are not required.

\begin{theorem}[Gaussian-core local convergence]
\label{thm:local-gaussian-core}
Let \(0<\gamma_0\leq\gstar\), and define
\begin{equation}
  \rho_{\mathrm{core}}(\delta)
  :=
  2\gamma_0-2+2\ell_\delta
  -\Kng(\delta)\mathfrak b(\delta)\sqrt\delta.
  \label{eq:local-core-rate}
\end{equation}
If \(\rho_{\mathrm{core}}(\delta)>0\) and
\(\Delta(a_0)\leq\delta\), then the Fisher--Rao flow remains in
\(\mathcal U_\delta\) and
\begin{equation}
  \norm{a_t-a_\star}_\star^2
  \leq
  e^{-\rho_{\mathrm{core}}(\delta)t}
  \norm{a_0-a_\star}_\star^2.
  \label{eq:local-core-contraction}
\end{equation}
With
\(\gamma_0=\max\{\alphastar,(4+\Gamma)^{-1}\}\), the limiting
squared-norm rate as \(\delta\downarrow0\) is at least
\(2\gamma_0\).
\end{theorem}

\begin{proof}
Energy dissipation makes \(\mathcal U_\delta\) forward invariant.
Equation \eqref{eq:local-gaussian-dissipation} contributes
\(-2\ell_\delta\norm{\xi}_\star^2\) to the derivative of the squared
norm.  The linear part of \(H\) contributes at most
\(2(1-\gamma_0)\norm{\xi}_\star^2\).  Taylor's theorem on
\(\mathcal H_\delta\), together with
\eqref{eq:local-norm-conversion}, bounds the remaining contribution by
\(\Kng(\delta)\mathfrak b(\delta)\sqrt\delta
\norm{\xi}_\star^2\).  The result follows from Gr\"onwall's inequality.
\end{proof}

\begin{corollary}[Exact Gaussian region]
\label{cor:local-exact-gaussian-region}
For a Gaussian target, \(\gstar=1\), \(\Kng\equiv0\), and
\(\rho_{\mathrm{core}}(\delta)=2\ell_\delta\).  For every
\(\rho\in(0,2)\), the largest energy sublevel on which the flow has
squared-norm rate at least \(\rho\) is
\begin{equation}
  \Delta\leq
  \Delta_{\mathrm G}^{\sharp}(\rho)
  :=
  \frac12\left(
    \frac{\rho}{2}-1-\log\frac{\rho}{2}
  \right).
  \label{eq:local-exact-gaussian-threshold}
\end{equation}
This threshold is independent of dimension.
\end{corollary}

\begin{proof}
For the standard Gaussian, each covariance eigenvalue solves
\(\dot c=c(1-c)\).  The instantaneous squared-norm rate is at least
\(2\lambda_{\min}(C)\), with equality for a centered covariance perturbation
supported on a minimal eigendirection.  If
\(\Delta\leq\Delta_{\mathrm G}^{\sharp}(\rho)\), the Gaussian energy formula
forces every covariance eigenvalue to be at least \(\rho/2\), so the rate is
at least \(\rho\).  Conversely, setting one eigenvalue to
\(c<\rho/2\) gives rate \(2c<\rho\) at energy cost
\(\frac12(c-1-\log c)\).  Letting \(c\uparrow\rho/2\) shows that every
strictly larger energy sublevel contains such a state, proving sharpness.
\end{proof}

\subsection{Discrete local contractions and entry}
\label{sec:local-discrete}

For the KL map, the Gaussian covariance eigenvalue and mean factors are
\[
  c^+=\frac{(1+h)c}{1+hc},
  \qquad
  r_C(c)=(1+hc)^{-2},
  \qquad
  r_m(c)=(1-hc)^2.
\]
Hence, for \(q\in(0,1)\), the exact invariant interval producing contraction
factor \(q\) is
\begin{equation}
  I_h^{\KL}(q)
  =
  \left[
    \frac{q^{-1/2}-1}{h},
    \frac{1+\sqrt q}{h}
  \right],
  \label{eq:local-kl-invariant-interval}
\end{equation}
provided \(0<h<2\) and
\(q\geq\max\{(1+h)^{-2},(1-h)^2\}\).
Indeed, the two scalar inequalities \(r_C(c),r_m(c)\leq q\) give exactly
the endpoints in \eqref{eq:local-kl-invariant-interval}, and the condition on
\(q\) is precisely what makes that interval contain \(1\).  Since \(c^+\)
lies between \(c\) and \(1\), the interval is forward invariant.  Hence
\[
  \operatorname{spec}(C_n)\subseteq I_h^{\KL}(q)
  \quad\Longrightarrow\quad
  \norm{a_{n+k}-a_\star}_{\star,h}^2
  \leq q^k\norm{a_n-a_\star}_{\star,h}^2,\qquad k\geq0.
\]

For non-Gaussian targets, let \(T_\bullet\) denote either deterministic map
and use the norm in \Cref{cor:local-linear-maps}.  Define
\begin{equation}
  K_{\mathrm{disc},\bullet}(\delta)
  :=
  h^{-1}
  \sup_{\substack{a\in\mathcal H_\delta\\
                   \norm{\xi}_{\star,\bullet}=1}}
  \norm{D^2T_\bullet(a)[\xi,\xi]}_{\star,\bullet}.
  \label{eq:local-discrete-modulus}
\end{equation}
Here
\(\norm{\cdot}_{\star,\Riem}=\norm{\cdot}_\star\),
\(\norm{\cdot}_{\star,\KL}=\norm{\cdot}_{\star,h}\), and
\begin{equation}
  \mathfrak b_{\Riem}(\delta):=\mathfrak b(\delta),
  \qquad
  \mathfrak b_{\KL}(\delta):=\mathfrak b_h(\delta).
  \label{eq:local-discrete-b-factors}
\end{equation}
The supremum is finite because Gaussian score differentiation makes the two
maps twice continuously differentiable on a neighborhood of the compact,
spectrally bounded hull; compactness by itself would not be enough.  For
fixed \(h\), the same continuity gives
\(K_{\mathrm{disc},\bullet}(\delta)
\to K_{\mathrm{disc},\bullet}(0)<\infty\) as
\(\delta\downarrow0\), where the value at zero is the same supremum with
\(a=a_\star\).  Since
\(\mathfrak b_\bullet(\delta)\sqrt\delta\to0\), the radius condition below
therefore holds for every sufficiently small positive \(\delta\).

\begin{theorem}[Local deterministic maps]
\label{thm:local-discrete-contraction}
Suppose
\[
  0<h\leq
  \min\!\left\{
    \frac{2}{4+\Gamma},
    \frac{1}{\betastar\max\{U_\delta,\alphastar^{-1}\}}
  \right\}
\]
and
\[
  K_{\mathrm{disc},\bullet}(\delta)
  \mathfrak b_\bullet(\delta)\sqrt\delta
  \leq\gamma_\bullet,
  \qquad
  \gamma_{\Riem}:=\frac1{4+\Gamma},\quad
  \gamma_{\KL}:=\frac1{4+2h+\Gamma}.
\]
If \(\Delta(a_0)\leq\delta\), then the iterates remain in
\(\mathcal U_\delta\) and
\begin{equation}
  \norm{a_n-a_\star}_{\star,\bullet}
  \leq
  \left(1-\frac{h\gamma_\bullet}{2}\right)^n
  \norm{a_0-a_\star}_{\star,\bullet}.
  \label{eq:local-discrete-contraction}
\end{equation}
\end{theorem}

\begin{proof}
Put
\(\overline\Lambda_\delta
:=\max\{U_\delta,\alphastar^{-1}\}\).
The stepsize restriction and
\Cref{prop:global:one-step-descent}, applied in optimizer-whitened
coordinates with curvature bound \(\betastar\) and covariance upper barrier
\(\overline\Lambda_\delta\), imply
\(\calE(T_\bullet(a))\leq\calE(a)\) for every
\(a\in\mathcal U_\delta\).  Thus \(\mathcal U_\delta\) is invariant by
induction.

Write \(\xi_n=a_n-a_\star\) and
\(M_\bullet=DT_\bullet(a_\star)\).  Since
\(T_\bullet(a_\star)=a_\star\), the integral Taylor formula gives
\[
  a_{n+1}-a_\star
  =M_\bullet\xi_n
   +\int_0^1(1-s)
      D^2T_\bullet(a_\star+s\xi_n)[\xi_n,\xi_n]\,\dd s.
\]
The whole segment lies in \(\mathcal H_\delta\).  By
\Cref{cor:local-linear-maps}, the linear term has norm at most
\((1-h\gamma_\bullet)\norm{\xi_n}_{\star,\bullet}\), whereas
\eqref{eq:local-discrete-modulus}, the radius hypothesis, and the
energy-to-norm conversion with \(\mathfrak b_\bullet\) bound the remainder by
\[
  \frac h2K_{\mathrm{disc},\bullet}(\delta)
  \norm{\xi_n}_{\star,\bullet}^2
  \leq
  \frac{h\gamma_\bullet}{2}
  \norm{\xi_n}_{\star,\bullet}.
\]
Adding the two estimates and iterating proves
\eqref{eq:local-discrete-contraction}.
\end{proof}

\begin{remark}[What is explicit]
\label{rem:local-discrete-explicitness}
The spectral gap, the stepsize restriction, the energy-to-norm conversion,
and the exact Gaussian intervals are explicit.  For a general target, the
certified radius is constructive through
\eqref{eq:local-discrete-modulus}, but the current theory does not provide a
closed-form bound for that supremum in the target parameters.  We therefore
do not call the general discrete radius fully explicit.
\end{remark}

An online switch need not know \(\calE(a_\star)\).  In original coordinates,
set
\begin{equation}
  \mathcal R_{\BW}(a)^2
  :=
  \norm{g(a)}^2+
  \Tr\!\left[
    (C^{-1}-A(a))C(C^{-1}-A(a))
  \right],
  \label{eq:local-bw-residual}
\end{equation}
where \(g(a)=-\E_{\N(m,C)}\nabla V\) and
\(A(a)=\E_{\N(m,C)}\nabla^2V\).
The Bures--Wasserstein PL inequality gives
\begin{equation}
  \mathcal R_{\BW}(a)^2\geq2\alpha\Delta(a).
  \label{eq:local-bw-gate}
\end{equation}
Thus
\(\mathcal R_{\BW}(a_n)^2\leq2\alpha\delta_{\mathrm{loc}}\)
certifies entry into any previously validated local level
\(\delta_{\mathrm{loc}}\).  This is an exact-expectation gate.  A Monte Carlo
implementation requires a concentration margin, and the stochastic theory
still needs its separate containment argument.

\subsection{A sharp growing-dimensional obstruction}
\label{sec:local-ridge}

The universal lower bound in
\Cref{thm:local-spectral-sandwich} cannot be replaced by a
dimension-free logarithmic rate under a globally Lipschitz Hessian assumption
alone.

\begin{theorem}[Smooth convex-ridge obstruction]
\label{thm:local-ridge-obstruction}
There are numerical constants \(c,C>0\) and a sequence of smooth strongly
convex potentials \(V_K:\R^{d_K}\to\R\), \(K\to\infty\), such that the
standard Gaussian is the optimizing Gaussian,
\[
  \E\nabla V_K(Z)=0,\qquad
  \E\nabla^2V_K(Z)=\Id,
\]
and the Hessians have a common global Lipschitz constant.  Moreover,
\begin{equation}
  d_K=\Theta(K^4),\qquad
  \kappastar=\Theta(K^2),\qquad
  \norm{\mathcal T}^2=\Theta(K),
  \qquad
  \gstar=\Theta(K^{-1}).
  \label{eq:local-ridge-scaling}
\end{equation}
Equivalently,
\[
  d_K=\Theta(\kappastar^2),\qquad
  \norm{\mathcal T}^2=\Theta(\kappastar^{1/2}),
  \qquad
  \gstar=\Theta(\kappastar^{-1/2}).
\]
\end{theorem}

\begin{proof}[Proof sketch]
The full construction and proof, given in
\Cref{app:local-ridge-construction}, use a convex ridge with one
distinguished coordinate and a weak tilt toward a radial block of
\(d_K-1\) coordinates.  Its averaged Hessian is tuned to be the identity.
A distributional Codazzi identity exposes a collective radial fluctuation
that gives
\(\norm{\mathcal T[X]}^2\geq K/2-O(1)\) for an explicit unit
Frobenius matrix \(X\).  The polarized second-Hermite identity makes the
packed generator self-adjoint.  The Rayleigh vector
\[
  Y=\frac{X}{\sqrt2},
  \qquad
  v=-\frac{\mathcal T[X]}2
\]
has quotient
\[
  \frac{2+\inner{X}{\mathcal H[X]}_{\Frob}
    -\norm{\mathcal T[X]}^2}
       {2+\norm{\mathcal T[X]}^2}
  =O(K^{-1}),
\]
which proves \(\gstar\leq C/K\).  On this family
\(\Gamma=\betastar-1=\Theta(K)\); the lower bound in
\Cref{thm:local-spectral-sandwich} therefore gives
\(\gstar\geq c/K\).  Mollifying the one-dimensional ridge profile and
retuning its scalar tilt preserve the exact whitening, the scaling, and the
common Hessian-Lipschitz bound.
\end{proof}

\begin{remark}[Scope of sharpness]
\label{rem:local-ridge-scope}
The construction saturates the \(\betastar-1\) branch of \(\Gamma\), not the
\(\sqrt d\log\kappastar\) branch.  Its dimension grows as
\(\Theta(\kappastar^2)\).  It rules out a dimension-uniform
\((1+\log\kappastar)^{-1}\) local rate, but it neither supplies a
fixed-dimensional counterexample nor determines the optimal dimension
dependence.
\end{remark}

\subsection{Global-to-local composition}
\label{sec:local-global-to-local}

The preceding results now compose without identifying the global energy gap
with the local parameter error.  Put
\[
  \underline\gamma_\star
  :=\max\!\left\{\alphastar,\frac1{4+\Gamma}\right\}.
\]

\begin{theorem}[Deterministic three-stage complexity]
\label{thm:local-three-stage}
Assume the global curvature condition
\eqref{eq:global:curvature-assumption} and the initial covariance band
\eqref{eq:global:initial-covariance-band}; in particular,
\(\Delta_0\), \(\lambda_0\), and \(\lambda_{\max}\) have the meanings fixed
in \Cref{sec:global}, while \(\lambda_{0,\star}\) is defined in
\eqref{eq:global:intrinsic-initial-floor}.
Let \(\delta_{\mathrm c}>0\) satisfy the hypotheses of
\Cref{thm:local-gaussian-core} with
\(\gamma_0=\underline\gamma_\star\) and
\(\rho_{\mathrm{core}}(\delta_{\mathrm c})
\geq\underline\gamma_\star\).  Let
\(\delta_{\Riem}\) and \(\delta_{\KL}\) satisfy the hypotheses of
\Cref{thm:local-discrete-contraction} for the indicated schemes.  Such
levels exist: \(\rho_{\mathrm{core}}(\delta)\to
2\underline\gamma_\star\), and, for each fixed admissible \(h\), the
discrete radius condition holds at every sufficiently small positive level.

For the continuous flow, squared local error at most \(\varepsilon\) is
reached by time
\begin{equation}
  T_{\mathrm c}(\varepsilon)
  =O\!\left(
  \begin{aligned}
    &\logp\frac1{\betastar\lambda_{0,\star}}
    +\kappastar\logp\frac{\Delta_0}{\delta_{\mathrm c}}\\[-1mm]
    &\qquad +(4+\Gamma)
      \logp\frac{\mathfrak b(\delta_{\mathrm c})^2
                       \delta_{\mathrm c}}{\varepsilon}
  \end{aligned}
  \right),
  \label{eq:local-three-stage-continuous}
\end{equation}
where the first line may be replaced by the smaller original-coordinate
global bound from \Cref{thm:global:continuous-upper}.

For the Riemannian map, suppose
\[
  0<h\leq\min\!\left\{
    \frac1{2\beta\lambda_{\max}},
    \frac2{4+\Gamma},
    \frac1{\betastar
      \max\{U_{\delta_{\Riem}},\alphastar^{-1}\}}
  \right\}.
\]
Then \(\norm{a_N-a_\star}_\star^2\leq\varepsilon\) after
\begin{equation}
  N_{\Riem}(\varepsilon)
  =O\!\left(
    \frac1h\logp\frac1{\beta\lambda_0}
    +\frac\kappa h\logp\frac{\Delta_0}{\delta_{\Riem}}
    +\frac{4+\Gamma}{h}
       \logp\frac{\mathfrak b(\delta_{\Riem})^2
                        \delta_{\Riem}}{\varepsilon}
  \right).
  \label{eq:local-three-stage-riemannian}
\end{equation}

For the KL map, suppose
\[
  0<h\leq\min\!\left\{
    \frac1{\beta\lambda_{\max}},
    \frac2{4+\Gamma},
    \frac1{\betastar
      \max\{U_{\delta_{\KL}},\alphastar^{-1}\}}
  \right\}.
\]
Then \(\norm{a_N-a_\star}_{\star,h}^2\leq\varepsilon\) after
\begin{equation}
  N_{\KL}(\varepsilon)
  =O\!\left(
    \frac1h\logp\frac1{\beta\lambda_0}
    +\frac\kappa h\logp\frac{\Delta_0}{\delta_{\KL}}
    +\frac{4+2h+\Gamma}{h}
       \logp\frac{\mathfrak b_h(\delta_{\KL})^2
                        \delta_{\KL}}{\varepsilon}
  \right).
  \label{eq:local-three-stage-kl}
\end{equation}
The hidden constants are numerical once the certified entry levels are
fixed independently of \(\varepsilon\).
\end{theorem}

\begin{proof}
Apply \Cref{thm:global:continuous-upper} until the continuous trajectory
reaches \(\mathcal U_{\delta_{\mathrm c}}\).  At entry,
\eqref{eq:local-norm-conversion} bounds the squared local norm by
\(\mathfrak b(\delta_{\mathrm c})^2\delta_{\mathrm c}\), and
\Cref{thm:local-gaussian-core} contracts it at rate at least
\(\underline\gamma_\star\geq(4+\Gamma)^{-1}\).  This proves
\eqref{eq:local-three-stage-continuous}.

For either discrete scheme, \Cref{thm:global:discrete-upper} supplies the
covariance-bootstrap and global-localization terms.  The energy-to-norm
conversion at entry is \(\mathfrak b_\bullet(\delta_\bullet)^2
\delta_\bullet\).  Restarting
\Cref{thm:local-discrete-contraction} and using
\(-\log(1-x)\geq x\) gives the final term in
\eqref{eq:local-three-stage-riemannian} or
\eqref{eq:local-three-stage-kl}.
\end{proof}

The first two stages above are measured in energy, while the last is measured
in an optimizer-frozen parameter norm.  Keeping that change of error metric
explicit prevents a local spectral statement from being mistaken for a
global KL contraction.
